\ssr{ПРИЛОЖЕНИЕ Б}

\begingroup
% Настройка заголовка: справа, без переноса
\captionsetup[longtable]{singlelinecheck=false, justification=raggedleft}

\begin{longtable}{|p{0.2\textwidth}|p{0.08\textwidth}|p{0.63\textwidth}|}

% --- ЗАГОЛОВОК НА ПЕРВОЙ СТРАНИЦЕ ---
\caption{Таблица обнаружений реализации DeepSeek} \label{tab:deepseek_errors} \\
\hline
\textbf{Файл} & \textbf{Кол-во} & \textbf{Примеры ошибок (токен и расположение)} \\ \hline
\endfirsthead

% --- ЗАГОЛОВОК НА СЛЕДУЮЩИХ СТРАНИЦАХ ---
\multicolumn{3}{c}{\small\textit{Продолжение таблицы \ref{tab:deepseek_errors}}} \\
\hline
\textbf{Файл} & \textbf{Кол-во} & \textbf{Примеры ошибок (токен и расположение)} \\ \hline
\endhead

% --- НИЗ ---
\hline
\endfoot
\hline
\endlastfoot

% --- ДАННЫЕ ---

% main (2).pdf
main (2).pdf & 281 & \seqsplit{'«Московский'} (страница 1, строка 4) \\
& & \seqsplit{'университет)»'} (страница 1, строка 6) \\
& & \seqsplit{'«Информатика'} (страница 1, строка 8) \\
& & \seqsplit{'«Программное'} (страница 1, строка 9) \\
& & \seqsplit{'«Метод'} (страница 1, строка 13) \\
& & \seqsplit{'месте»'} (страница 1, строка 14) \\
& & \seqsplit{'•провести'} (страница 2, строка 9) \\
& & \seqsplit{'усталости\textasciitilde{}'} (страница 2, строка 9) \\
& & \seqsplit{'•провести'} (страница 2, строка 10) \\
& & \seqsplit{'места\textasciitilde{}'} (страница 2, строка 11) \\
& & \seqsplit{'•определить'} (страница 2, строка 12) \\
& & \seqsplit{'характеристик\textasciitilde{}'} (страница 2, строка 12) \\
& & \seqsplit{'•разработать'} (страница 2, строка 13) \\
& & \seqsplit{'оператора\textasciitilde{}'} (страница 2, строка 13) \\
& & \seqsplit{'•реализовать'} (страница 2, строка 14) \\
& & \seqsplit{'•нарушение'} (страница 7, строка 4) \\
& & \seqsplit{'•неудовлетворенность'} (страница 7, строка 5) \\
& & \seqsplit{'•проблемы'} (страница 7, строка 6) \\
& & \seqsplit{'•снижение'} (страница 7, строка 7) \\
& & \seqsplit{'•увеличение'} (страница 7, строка 8) \\
& & \textit{... и еще 261 шт.} \\ \hline

% _09.pdf
\_09.pdf & 15 & \seqsplit{'де\$\textbackslash{}sigma\$'} (страница 5, строка 8) \\
& & \seqsplit{'де\$\textbackslash{}alpha\$'} (страница 5, строка 12) \\
& & \seqsplit{'деP='} (страница 5, строка 48) \\
& & \seqsplit{'деP='} (страница 6, строка 3) \\
& & \seqsplit{'значение\$\textbackslash{}alpha\$'} (страница 6, строка 5) \\
& & \seqsplit{'значение\$\textbackslash{}alpha\$'} (страница 6, строка 6) \\
& & \seqsplit{'де\$\textbackslash{}sigma\$'} (страница 6, строка 27) \\
& & \seqsplit{'остатков\$\textbackslash{}varepsilon\$'} (страница 6, строка 39) \\
& & \seqsplit{'параметра\$\textbackslash{}alpha\$'} (страница 6, строка 42) \\
& & \seqsplit{'де\$\textbackslash{}sigma\$'} (страница 7, строка 10) \\
& & \seqsplit{'де\$\textbackslash{}gamma\$'} (страница 7, строка 14) \\
& & \seqsplit{'де\$\textbackslash{}gamma\$'} (страница 7, строка 18) \\
& & \seqsplit{'де\$\textbackslash{}sigma\$'} (страница 7, строка 23) \\
& & \seqsplit{'деh'} (страница 8, строка 10) \\
& & \seqsplit{'де\$\textbackslash{}sigma\$'} (страница 11, строка 4) \\ \hline

% ВКР Селез.pdf
ВКР Селез.pdf & 170 & \seqsplit{'«молдавского»'} (страница 3, строка 17) \\
& & \seqsplit{'«румынский»6,'} (страница 3, строка 18) \\
& & \seqsplit{'[и'} (страница 3, строка 28) \\
& & \seqsplit{'издание].'} (страница 3, строка 30) \\
& & \seqsplit{'«мягкой'} (страница 3, строка 35) \\
& & \seqsplit{'силы»'} (страница 3, строка 35) \\
& & \seqsplit{'Молдавии11.'} (страница 3, строка 37) \\
& & \seqsplit{'ЕС.'} (страница 4, строка 4) \\
& & \seqsplit{'юга/востока'} (страница 4, строка 4) \\
& & \seqsplit{'севера/запада'} (страница 4, строка 4) \\
& & \seqsplit{'(Пророссийская/Проевропейская/Любая)'} (страница 5, строка 17) \\
& & \seqsplit{'(Юг/Запад/Восток/Центр/Приднестро-'} (страница 5, строка 18) \\
& & \seqsplit{'вье/Кишинёв/Любой)'} (страница 5, строка 19) \\
& & \seqsplit{'экономика/Фармацевтика/-'} (страница 5, строка 27) \\
& & \seqsplit{'хозяйство/Любая)'} (страница 5, строка 28) \\
& & \seqsplit{'(ПАС/Партия'} (страница 5, строка 29) \\
& & \seqsplit{'«Шор»/Любая)'} (страница 5, строка 29) \\
& & \seqsplit{'«Шор».'} (страница 5, строка 32) \\
& & \seqsplit{'(Россия/ЕС/Польша/Великобритания/Франция/Израиль-'} (страница 6, строка 4) \\
& & \seqsplit{'/Греция/Любая)'} (страница 6, строка 5) \\
& & \textit{... и еще 150 шт.} \\ \hline

% _08.pdf
\_08.pdf & 2 & \seqsplit{'(’причастие'} (страница 4, строка 29) \\
& & \seqsplit{'(’причастие'} (страница 4, строка 31) \\ \hline

% _03.pdf
\_03.pdf & 39 & \seqsplit{'AиBразмерностей'} (страница 3, строка 5) \\
& & \seqsplit{'m\$\textbackslash{}times\$nиn\$\textbackslash{}times\$qсоответственно.'} (страница 3, строка 5) \\
& & \seqsplit{'Cразмерности'} (страница 3, строка 6) \\
& & \seqsplit{'m\$\textbackslash{}times\$qвычисляются'} (страница 3, строка 6) \\
& & \seqsplit{'формуле:'} (страница 3, строка 6) \\
& & \seqsplit{'умноженияAнаBтребуется'} (страница 3, строка 10) \\
& & \seqsplit{'выполнитьm\$\textbackslash{}cdot\$n\$\textbackslash{}cdot\$qопераций'} (страница 3, строка 10) \\
& & \seqsplit{'сложения.В'} (страница 3, строка 11) \\
& & \seqsplit{'O(n3)'} (страница 3, строка 12) \\
& & \seqsplit{'O(nlog27)'} (страница 3, строка 13) \\
& & \seqsplit{'\$\textbackslash{}approx\$O(n2.81)'} (страница 3, строка 13) \\
& & \seqsplit{'[1].Для'} (страница 3, строка 14) \\
& & \seqsplit{'алгоритма.'} (страница 3, строка 14) \\
& & \seqsplit{'матрицAиB,'} (страница 3, строка 15) \\
& & \seqsplit{'размераn\$\textbackslash{}times\$n,'} (страница 3, строка 16) \\
& & \seqsplit{'n=2k.'} (страница 3, строка 16) \\
& & \seqsplit{'матрицуAи'} (страница 3, строка 17) \\
& & \seqsplit{'Bразбиваем'} (страница 3, строка 17) \\
& & \seqsplit{'подматрицыA11,A12,A21,A22'} (страница 3, строка 18) \\
& & \seqsplit{'B11,B12,B21,B22'} (страница 3, строка 18) \\
& & \textit{... и еще 19 шт.} \\ \hline

% _02.pdf
\_02.pdf & 3798 & \seqsplit{'гдеX-'} (страница 4, строка 10) \\
& & \seqsplit{'i-е'} (страница 4, строка 11) \\
& & \seqsplit{'где\$\textbackslash{}mu\$Xи\$\textbackslash{}mu\$Y-'} (страница 4, строка 22) \\
& & \seqsplit{'значенияXиY'} (страница 4, строка 23) \\
& & \seqsplit{'\$\textbackslash{}sigma\$Xи\$\textbackslash{}sigma\$Y-'} (страница 4, строка 24) \\
& & \seqsplit{'отклоненияXиY'} (страница 4, строка 25) \\
& & \seqsplit{'E[(X'} (страница 4, строка 26) \\
& & \seqsplit{'\$\textbackslash{}mu\$X)(Y'} (страница 4, строка 26) \\
& & \seqsplit{'\$\textbackslash{}mu\$Y)]-'} (страница 4, строка 26) \\
& & \seqsplit{'ковариацияXиY'} (страница 4, строка 27) \\
& & \seqsplit{'rпринимает'} (страница 4, строка 28) \\
& & \seqsplit{'-1до1.'} (страница 4, строка 28) \\
& & \seqsplit{'r='} (страница 4, строка 29) \\
& & \seqsplit{'r='} (страница 4, строка 31) \\
& & \seqsplit{'-1,'} (страница 4, строка 31) \\
& & \seqsplit{'r=0,'} (страница 4, строка 32) \\
& & \seqsplit{'Спирмена[1].'} (страница 4, строка 33) \\
& & \seqsplit{'Коэффициент\$\textbackslash{}rho\$вычисляется'} (страница 4, строка 34) \\
& & \seqsplit{'формуле:'} (страница 4, строка 34) \\
& & \seqsplit{'di-'} (страница 4, строка 37) \\
& & \textit{... и еще 3778 шт.} \\ \hline

% _00.pdf
\_00.pdf & 1 & \seqsplit{'input("Введите'} (страница 7, строка 35) \\ \hline

% _01.pdf
\_01.pdf & 8 & \seqsplit{'«Заполнить\_массив\_операций'} (страница 6, строка 2) \\
& & \seqsplit{'«Изделие'} (страница 6, строка 3) \\
& & \seqsplit{'«Изделие'} (страница 7, строка 2) \\
& & \seqsplit{'«Заполнить'} (страница 7, строка 3) \\
& & \seqsplit{'массив'} (страница 7, строка 3) \\
& & \seqsplit{'операций»'} (страница 7, строка 4) \\
& & \seqsplit{'«Удаление'} (страница 8, строка 2) \\
& & \seqsplit{'изделия»'} (страница 8, строка 2) \\ \hline

% _65-1.pdf
\_65-1.pdf & 29 & \seqsplit{'иB'} (страница 4, строка 10) \\
& & \seqsplit{'иB'} (страница 5, строка 5) \\
& & \seqsplit{'деani,'} (страница 5, строка 27) \\
& & \seqsplit{'дoёмких'} (страница 5, строка 30) \\
& & \seqsplit{'дoёмк'} (страница 11, строка 8) \\
& & \seqsplit{'дoёмк'} (страница 11, строка 10) \\
& & \seqsplit{'словия+\$\textbackslash{}begin\{cases\}\$'} (страница 11, строка 12) \\
& & \seqsplit{'дoёмк'} (страница 11, строка 16) \\
& & \seqsplit{'сравнения+N(f'} (страница 11, строка 17) \\
& & \seqsplit{'тела+f'} (страница 11, строка 17) \\
& & \seqsplit{'дoёмк'} (страница 11, строка 18) \\
& & \seqsplit{'функции/возврат'} (страница 11, строка 18) \\
& & \seqsplit{'дoёмк'} (страница 12, строка 11) \\
& & \seqsplit{'дoёмк'} (страница 12, строка 12) \\
& & \seqsplit{'дoёмк'} (страница 12, строка 13) \\
& & \seqsplit{'дoёмк'} (страница 12, строка 19) \\
& & \seqsplit{'дoёмк'} (страница 13, строка 1) \\
& & \seqsplit{'дoёмк'} (страница 13, строка 4) \\
& & \seqsplit{'дoёмк'} (страница 13, строка 7) \\
& & \seqsplit{'дoёмк'} (страница 13, строка 11) \\
& & \textit{... и еще 9 шт.} \\ \hline

% _04.pdf
\_04.pdf & 329 & \seqsplit{'Cценарий'} (страница 3, строка 22) \\
& & \seqsplit{'работеThe'} (страница 7, строка 2) \\
& & \seqsplit{'---Аффилиативный'} (страница 7, строка 6) \\
& & \seqsplit{'---Барьеры'} (страница 7, строка 7) \\
& & \seqsplit{'---Бизнес-ангел'} (страница 7, строка 8) \\
& & \seqsplit{'---Бизнес-инкубатор'} (страница 7, строка 9) \\
& & \seqsplit{'---Бизнес-модель'} (страница 7, строка 10) \\
& & \seqsplit{'---Бизнес-план'} (страница 7, строка 11) \\
& & \seqsplit{'---Бюджетирование'} (страница 7, строка 12) \\
& & \seqsplit{'---Венчурный'} (страница 7, строка 13) \\
& & \seqsplit{'---Грант'} (страница 7, строка 14) \\
& & \seqsplit{'---Диверсификация'} (страница 7, строка 15) \\
& & \seqsplit{'---Дивиденды'} (страница 7, строка 16) \\
& & \seqsplit{'---Инвестиции'} (страница 7, строка 17) \\
& & \seqsplit{'---Инновация'} (страница 7, строка 18) \\
& & \seqsplit{'---Капитализация'} (страница 7, строка 19) \\
& & \seqsplit{'---Конкурентное'} (страница 7, строка 20) \\
& & \seqsplit{'---Краудфандинг'} (страница 7, строка 21) \\
& & \seqsplit{'---Ликвидность'} (страница 7, строка 22) \\
& & \seqsplit{'---Маркетинг'} (страница 7, строка 23) \\
& & \textit{... и еще 309 шт.} \\ \hline

% _10.pdf
\_10.pdf & 68 & \seqsplit{'анализ---'} (страница 3, строка 5) \\
& & \seqsplit{'помощьюDefinite'} (страница 3, строка 9) \\
& & \seqsplit{'nи'} (страница 3, строка 20) \\
& & \seqsplit{'n.'} (страница 3, строка 20) \\
& & \seqsplit{'интеграл\$\textbackslash{}int\$f(x)dx.'} (страница 3, строка 22) \\
& & \seqsplit{'методы:'} (страница 3, строка 26) \\
& & \seqsplit{'прямоугольников;'} (страница 3, строка 27) \\
& & \seqsplit{'трапеций;'} (страница 3, строка 28) \\
& & \seqsplit{'Симпсона.'} (страница 3, строка 29) \\
& & \seqsplit{'метод:'} (страница 3, строка 30) \\
& & \seqsplit{'формула Ньютона-Лейбница.'} (страница 3, строка 31) \\
& & \seqsplit{'Python.'} (страница 3, строка 33) \\
& & \seqsplit{'Библиотеки:'} (страница 3, строка 33) \\
& & \seqsplit{'numpy,'} (страница 3, строка 34) \\
& & \seqsplit{'scipy,'} (страница 3, строка 35) \\
& & \seqsplit{'matplotlib.'} (страница 3, строка 36) \\
& & \seqsplit{'Результаты:'} (страница 3, строка 37) \\
& & \seqsplit{'сравнение'} (страница 3, строка 37) \\
& & \seqsplit{'точности.'} (страница 3, строка 37) \\
& & \seqsplit{'Выводы:'} (страница 3, строка 38) \\
& & \textit{... и еще 48 шт.} \\ \hline

% _06.pdf
\_06.pdf & 4 & \seqsplit{'чaс'} (страница 2, строка 8) \\
& & \seqsplit{'«белого'} (страница 9, строка 9) \\
& & \seqsplit{'«белого»'} (страница 10, строка 12) \\
& & \seqsplit{'ящику»'} (страница 10, строка 12) \\ \hline

% _07.pdf
\_07.pdf & 24 & \seqsplit{'«сущность-связь»,'} (страница 2, строка 5) \\
& & \seqsplit{'«Сырые»'} (страница 3, строка 6) \\
& & \seqsplit{'«Шор»):'} (страница 3, строка 19) \\
& & \seqsplit{'«Шор»'} (страница 3, строка 22) \\
& & \seqsplit{'«Шор»,'} (страница 3, строка 38) \\
& & \seqsplit{'«Шор»'} (страница 4, строка 4) \\
& & \seqsplit{'«Шор»'} (страница 4, строка 23) \\
& & \seqsplit{'«Шор»'} (страница 5, строка 12) \\
& & \seqsplit{'«Шор»'} (страница 5, строка 18) \\
& & \seqsplit{'«Шор»'} (страница 5, строка 21) \\
& & \seqsplit{'«Шор»'} (страница 5, строка 22) \\
& & \seqsplit{'«Шор»'} (страница 5, строка 23) \\
& & \seqsplit{'«Шор»'} (страница 5, строка 24) \\
& & \seqsplit{'«Шор»'} (страница 5, строка 37) \\
& & \seqsplit{'«Шор»'} (страница 6, строка 18) \\
& & \seqsplit{'«Шор»'} (страница 6, строка 19) \\
& & \seqsplit{'«Шор»'} (страница 6, строка 20) \\
& & \seqsplit{'«Шор»'} (страница 6, строка 21) \\
& & \seqsplit{'«Шор»,'} (страница 7, строка 16) \\
& & \seqsplit{'«Шор».'} (страница 7, строка 17) \\
& & \textit{... и еще 4 шт.} \\ \hline

\end{longtable}
\endgroup



\begingroup
% Настройка заголовка: справа, без переноса
\captionsetup[longtable]{singlelinecheck=false, justification=raggedleft}

\begin{longtable}{|p{0.2\textwidth}|p{0.65\textwidth}|p{0.08\textwidth}|}

% --- ЗАГОЛОВОК НА ПЕРВОЙ СТРАНИЦЕ ---
\caption{Таблица обнаружений реализации Qwen-Coder} \label{tab:qwen_errors} \\
\hline
\textbf{Файл} & \textbf{Примеры ошибок (токен)} & \textbf{Стр.} \\ \hline
\endfirsthead

% --- ЗАГОЛОВОК НА СЛЕДУЮЩИХ СТРАНИЦАХ ---
\multicolumn{3}{c}{\small\textit{Продолжение таблицы \ref{tab:qwen_errors}}} \\
\hline
\textbf{Файл} & \textbf{Примеры ошибок (токен)} & \textbf{Стр.} \\ \hline
\endhead

% --- НИЗ ---
\hline
\endfoot
\hline
\endlastfoot

% --- ДАННЫЕ ---

% _02.pdf
\_02.pdf & \seqsplit{'гдеX-'} (строка 10) & 4 \\
& \seqsplit{'i-е'} (строка 11) & 4 \\
& \seqsplit{'XиY'} (строка 18) & 4 \\
& \seqsplit{'Xи'} (строка 22) & 4 \\
& \seqsplit{'Xи'} (строка 30) & 4 \\
& \seqsplit{'гдеdi-'} (строка 4) & 5 \\
& \seqsplit{'топ-N'} (строка 22) & 6 \\
& \seqsplit{'meanмежду'} (строка 30) & 23 \\ \hline

% _03.pdf
\_03.pdf & \seqsplit{'AиBразмерностей'} (строка 5) & 3 \\
& \seqsplit{'Cразмерности'} (строка 6) & 3 \\
& \seqsplit{'Cкак'} (строка 11) & 3 \\
& \seqsplit{'qэлементов'} (строка 14) & 3 \\
& \seqsplit{'nскалярное'} (строка 19) & 3 \\
& \seqsplit{'иv'} (строка 20) & 3 \\
& \seqsplit{'nвводят'} (строка 2) & 4 \\
& \seqsplit{'еслиnнечётно'} (строка 13) & 4 \\
& \seqsplit{'еслиnчётно'} (строка 14) & 4 \\
& \seqsplit{'fусловия'} (строка 2) & 8 \\
& \seqsplit{'fинициализации'} (строка 7) & 8 \\
& \seqsplit{'cреда'} (строка 4) & 10 \\
& \seqsplit{'иrandom'} (строка 19) & 17 \\ \hline

% _04.pdf
\_04.pdf & \seqsplit{'Cценарий'} (строка 22) & 3 \\
& \seqsplit{'работеThe'} (строка 2) & 7 \\
& \seqsplit{'cистема'} (строка 27) & 8 \\
& \seqsplit{'cистема'} (строка 29) & 8 \\
& \seqsplit{'AI-модель'} (строка 21) & 10 \\
& \seqsplit{'AI-алгоритмы'} (строка 5) & 11 \\
& \seqsplit{'AI-системы'} (строка 12) & 12 \\
& \seqsplit{'AI-генерируемый'} (строка 20) & 12 \\
& \seqsplit{'Python-скрипт'} (строка 4) & 13 \\
& \seqsplit{'GIFДа'} (строка 15) & 16 \\
& \seqsplit{'generatorГенерация'} (строка 27) & 16 \\
& \seqsplit{'MachineГенерация'} (строка 9) & 17 \\
& \seqsplit{'CaptionsГенерация'} (строка 15) & 17 \\
& \seqsplit{'CaptioningГенерация'} (строка 25) & 17 \\
& \seqsplit{'HumorГенерация'} (строка 40) & 17 \\
& \seqsplit{'AdaptabilityГенерация'} (строка 49) & 17 \\
& \seqsplit{'GeneratorГенерация'} (строка 54) & 17 \\
& \seqsplit{'cuptionsГенерация'} (строка 59) & 17 \\
& \seqsplit{'GeneratorГенерация'} (строка 8) & 18 \\
& \seqsplit{'GeneratorГенерация'} (строка 14) & 18 \\
& \seqsplit{'GeneratorГенерация'} (строка 19) & 18 \\
& \seqsplit{'GeneratorГенерация'} (строка 31) & 18 \\
& \seqsplit{'GeneratorГенерация'} (строка 49) & 18 \\
& \seqsplit{'GeneratorГенерация'} (строка 4) & 19 \\
& \seqsplit{'GeneratorГенерация'} (строка 13) & 19 \\
& \seqsplit{'GeneratorГенерация'} (строка 18) & 19 \\
& \seqsplit{'IDEF0-диаграмма'} (строка 5) & 22 \\
& \seqsplit{'IDEF0-диаграмма'} (строка 7) & 22 \\
& \seqsplit{'xиG'} (строка 25) & 24 \\
& \seqsplit{'aрктангенс'} (строка 43) & 24 \\
& \seqsplit{'пикселяg'} (строка 2) & 25 \\
& \seqsplit{'гдеs'} (строка 6) & 25 \\
& \seqsplit{'сивностиh'} (строка 15) & 25 \\
& \seqsplit{'гдеn'} (строка 17) & 25 \\
& \seqsplit{'ЕслиF'} (строка 24) & 25 \\
& \seqsplit{'картыP'} (строка 25) & 25 \\
& \seqsplit{'гдеHиW'} (строка 30) & 25 \\
& \seqsplit{'частиx'} (строка 34) & 25 \\
& \seqsplit{'iвычисляются'} (строка 1) & 26 \\
& \seqsplit{'частьюx'} (строка 2) & 26 \\
& \seqsplit{'признакиx'} (строка 16) & 26 \\
& \seqsplit{'классуC'} (строка 17) & 26 \\
& \seqsplit{'вероятностиP'} (строка 18) & 26 \\
& \seqsplit{'zkдля'} (строка 19) & 26 \\
& \seqsplit{'вx'} (строка 24) & 26 \\
& \seqsplit{'метокs'} (строка 27) & 26 \\
& \seqsplit{'вероятностьP'} (строка 28) & 26 \\
& \seqsplit{'энергиюE'} (строка 29) & 26 \\
& \seqsplit{'i-го'} (строка 10) & 27 \\
& \seqsplit{'i-го'} (строка 11) & 27 \\
& \seqsplit{'последовательностиSна'} (строка 9) & 28 \\
& \seqsplit{'кенамиt'} (строка 10) & 28 \\
& \seqsplit{'гдеc'} (строка 3) & 29 \\
& \seqsplit{'токенt'} (строка 6) & 29 \\
& \seqsplit{'обычноN'} (строка 7) & 29 \\
& \seqsplit{'токенуt'} (строка 8) & 29 \\
& \seqsplit{'ПустьT'} (строка 15) & 29 \\
& \seqsplit{'токенуv'} (строка 17) & 29 \\
& \seqsplit{'гдеd'} (строка 19) & 29 \\
& \seqsplit{'строкаW'} (строка 24) & 29 \\
& \seqsplit{'токеновT'} (строка 1) & 30 \\
& \seqsplit{'гдеe'} (строка 4) & 30 \\
& \seqsplit{'эмбеддинговX'} (строка 5) & 30 \\
& \seqsplit{'черезLидентичных'} (строка 6) & 30 \\
& \seqsplit{'токенx'} (строка 7) & 30 \\
& \seqsplit{'Значениеv'} (строка 8) & 30 \\
& \seqsplit{'гдеQ'} (строка 12) & 30 \\
& \seqsplit{'проекциямиWQ'} (строка 13) & 30 \\
& \seqsplit{'ходуz'} (строка 20) & 30 \\
& \seqsplit{'последнегоL-го'} (строка 23) & 30 \\
& \seqsplit{'векторовH'} (строка 24) & 30 \\
& \seqsplit{'последовательностиY'} (строка 27) & 30 \\
& \seqsplit{'токенаy'} (строка 28) & 30 \\
& \seqsplit{'промтаHи'} (строка 30) & 30 \\
& \seqsplit{'логитовl'} (строка 31) & 30 \\
& \seqsplit{'вероятностейP'} (строка 32) & 30 \\
& \seqsplit{'гдеl'} (строка 37) & 30 \\
& \seqsplit{'токенаy'} (строка 38) & 30 \\
& \seqsplit{'токенy'} (строка 2) & 31 \\
& \seqsplit{'генерацииy'} (строка 3) & 31 \\
& \seqsplit{'Cценарий'} (строка 6) & 31 \\
& \seqsplit{'IDEF0-диаграммы'} (строка 14) & 31 \\
& \seqsplit{'функцииget'} (строка 26) & 35 \\
& \seqsplit{'функцииgenerate'} (строка 24) & 37 \\
& \seqsplit{'женияMistral'} (строка 4) & 41 \\
& \seqsplit{'Моcква'} (строка 8) & 62 \\ \hline

% _06.pdf
\_06.pdf & \seqsplit{'IDEFО'} (строка 10) & 1 \\
& \seqsplit{'чaс'} (строка 8) & 2 \\
& \seqsplit{'IDEFО'} (строка 25) & 2 \\
& \seqsplit{'IDEFО'} (строка 7) & 10 \\ \hline

% _07.pdf
\_07.pdf & \seqsplit{'программированияPython'} (строка 3) & 7 \\
& \seqsplit{'классPoliticalRiskAnalyzer'} (строка 20) & 7 \\ \hline

% _09.pdf
\_09.pdf & \seqsplit{'деP'} (строка 48) & 5 \\
& \seqsplit{'ER-диаграмма'} (строка 28) & 6 \\
& \seqsplit{'ER-диаграмма'} (строка 1) & 7 \\
& \seqsplit{'CID-транзакций'} (строка 26) & 18 \\
& \seqsplit{'SQL-инъекций'} (строка 7) & 20 \\ \hline

% _10.pdf
\_10.pdf & \seqsplit{'помощьюDefinite'} (строка 9) & 3 \\
& \seqsplit{'nи'} (строка 20) & 3 \\
& \seqsplit{'всеB'} (строка 21) & 3 \\
& \seqsplit{'фактыresult'} (строка 26) & 3 \\
& \seqsplit{'Personконкретизируются'} (строка 4) & 4 \\
& \seqsplit{'4иverb'} (строка 9) & 4 \\
& \seqsplit{'DCG-правила'} (строка 13) & 4 \\
& \seqsplit{'intдаютfloat'} (строка 18) & 4 \\
& \seqsplit{'какint'} (строка 19) & 4 \\
& \seqsplit{'DCG-правило'} (строка 26) & 4 \\
& \seqsplit{'4иobject'} (строка 6) & 5 \\
& \seqsplit{'Personпри'} (строка 18) & 5 \\
& \seqsplit{'переменныеNumber'} (строка 27) & 5 \\
& \seqsplit{'рекурсивнаяDCG'} (строка 3) & 6 \\
& \seqsplit{'термовnoun'} (строка 7) & 7 \\
& \seqsplit{'floatи'} (строка 13) & 7 \\
& \seqsplit{'фактамиresult'} (строка 21) & 7 \\
& \seqsplit{'предикатаresult'} (строка 3) & 8 \\
& \seqsplit{'подставляетint'} (строка 6) & 8 \\
& \seqsplit{'DCG-правило'} (строка 21) & 9 \\
& \seqsplit{'какfloat'} (строка 26) & 9 \\
& \seqsplit{'типintпо'} (строка 27) & 9 \\
& \seqsplit{'иverb'} (строка 5) & 11 \\
& \seqsplit{'иverb'} (строка 7) & 11 \\
& \seqsplit{'иverb'} (строка 9) & 11 \\
& \seqsplit{'братverb'} (строка 15) & 11 \\
& \seqsplit{'еслиxнеизвестен'} (строка 1) & 12 \\
& \seqsplit{'обаint'} (строка 3) & 12 \\
& \seqsplit{'pluralиsingular'} (строка 17) & 12 \\
& \seqsplit{'системаможетбытьинтегрированассовременнымиNLP'} (строка 15) & 13 \\
& \seqsplit{'ОсобоевниманиепредполагаетсяуделитьинтеграцииссовременнымиNLP'} (строка 7) & 15 \\ \hline

% _65-1.pdf
\_65-1.pdf & \seqsplit{'иB'} (строка 10) & 4 \\
& \seqsplit{'иB'} (строка 5) & 5 \\
& \seqsplit{'деani'} (строка 27) & 5 \\
& \seqsplit{'дoёмких'} (строка 30) & 5 \\
& \seqsplit{'деM'} (строка 7) & 7 \\
& \seqsplit{'дoёмк'} (строка 8) & 11 \\
& \seqsplit{'дoёмк'} (строка 10) & 11 \\
& \seqsplit{'дoёмк'} (строка 16) & 11 \\
& \seqsplit{'дoёмк'} (строка 18) & 11 \\
& \seqsplit{'дoёмк'} (строка 11) & 12 \\
& \seqsplit{'дoёмк'} (строка 12) & 12 \\
& \seqsplit{'дoёмк'} (строка 13) & 12 \\
& \seqsplit{'иCOLS'} (строка 19) & 12 \\
& \seqsplit{'дoёмк'} (строка 1) & 13 \\
& \seqsplit{'дoёмк'} (строка 4) & 13 \\
& \seqsplit{'дoёмк'} (строка 7) & 13 \\
& \seqsplit{'дoёмк'} (строка 11) & 13 \\
& \seqsplit{'дoёмк'} (строка 16) & 13 \\
& \seqsplit{'дoёмк'} (строка 1) & 14 \\
& \seqsplit{'дoёмк'} (строка 3) & 14 \\
& \seqsplit{'дoёмк'} (строка 6) & 14 \\
& \seqsplit{'дoёмк'} (строка 9) & 14 \\
& \seqsplit{'деflast'} (строка 12) & 14 \\
& \seqsplit{'дoёмк'} (строка 17) & 14 \\
& \seqsplit{'дoёмких'} (строка 5) & 29 \\ \hline

% main (2).pdf
main (2).pdf & \seqsplit{'i-ый'} (строка 30) & 14 \\
& \seqsplit{'иy'} (строка 9) & 15 \\
& \seqsplit{'стиp'} (строка 1) & 20 \\
& \seqsplit{'c-средних'} (строка 4) & 30 \\
& \seqsplit{'k-средних'} (строка 26) & 30 \\
& \seqsplit{'гдеk'} (строка 27) & 30 \\
& \seqsplit{'c-средних'} (строка 6) & 34 \\
& \seqsplit{'HTTP-запросов'} (строка 3) & 40 \\
& \seqsplit{'YDVP-запрос'} (строка 8) & 40 \\
& \seqsplit{'YDVP-запросов'} (строка 12) & 40 \\
& \seqsplit{'YDVP-запросы'} (строка 16) & 40 \\
& \seqsplit{'YDVP-ответов'} (строка 27) & 40 \\
& \seqsplit{'c-средних'} (строка 10) & 46 \\
& \seqsplit{'POST-запросов'} (строка 26) & 47 \\
& \seqsplit{'POST-запросов'} (строка 29) & 47 \\
& \seqsplit{'мсPostgres'} (строка 28) & 49 \\
& \seqsplit{'c-средних'} (строка 32) & 50 \\
& \seqsplit{'c-средних'} (строка 19) & 54 \\
& \seqsplit{'BBНеактивен'} (строка 8) & 61 \\
& \seqsplit{'BBактивен'} (строка 38) & 63 \\
& \seqsplit{'BBнеактивен'} (строка 9) & 64 \\ \hline

% ВКР Селез.pdf
ВКР Селез.pdf & \seqsplit{'XIV-ХХI'} (строка 40) & 6 \\
& \seqsplit{'Oб'} (строка 27) & 7 \\
& \seqsplit{'Oб'} (строка 24) & 30 \\
& \seqsplit{'XIV-ХХI'} (строка 29) & 31 \\
& \seqsplit{'Oб'} (строка 13) & 68 \\
& \seqsplit{'XIV-ХХI'} (строка 29) & 70 \\ \hline

\end{longtable}
\endgroup


\begingroup
% Настройка заголовка: справа, без переноса
\captionsetup[longtable]{singlelinecheck=false, justification=raggedleft}

\begin{longtable}{|p{0.2\textwidth}|p{0.65\textwidth}|p{0.08\textwidth}|}

% --- ЗАГОЛОВОК НА ПЕРВОЙ СТРАНИЦЕ ---
\caption{Таблица обнаружений реализации ChatGPT 5.1} \label{tab:gpt_errors} \\
\hline
\textbf{Файл} & \textbf{Примеры ошибок (токен)} & \textbf{Стр.} \\ \hline
\endfirsthead

% --- ЗАГОЛОВОК НА СЛЕДУЮЩИХ СТРАНИЦАХ ---
\multicolumn{3}{c}{\small\textit{Продолжение таблицы \ref{tab:gpt_errors}}} \\
\hline
\textbf{Файл} & \textbf{Примеры ошибок (токен)} & \textbf{Стр.} \\ \hline
\endhead

% --- НИЗ ---
\hline
\endfoot
\hline
\endlastfoot

% --- ДАННЫЕ ---

% main (2).pdf
main (2).pdf & \seqsplit{ti—i-ый} (строка 30) & 14 \\
& \seqsplit{ωфункции} (строка 6) & 15 \\
& \seqsplit{x(t)иy(t)с} (строка 9) & 15 \\
& \seqsplit{стиp=} (строка 1) & 20 \\
& \seqsplit{c-средних.} (строка 4) & 30 \\
& \seqsplit{k-средних.} (строка 4) & 30 \\
& \seqsplit{k-средних} (строка 32) & 30 \\
& \seqsplit{c-средних} (строка 32) & 30 \\
& \seqsplit{гдеk—} (строка 33) & 30 \\
& \seqsplit{c-средних} (строка 6) & 34 \\
& \seqsplit{HTTP-запросов} (строка 3) & 40 \\
& \seqsplit{YDVP-запрос} (строка 8) & 40 \\
& \seqsplit{YDVP-ответ,} (строка 8) & 40 \\
& \seqsplit{YDVP-запросов.} (строка 12) & 40 \\
& \seqsplit{YDVP-запросы} (строка 16) & 40 \\
& \seqsplit{YDVP-ответов} (строка 27) & 40 \\
& \seqsplit{c-средних} (строка 10) & 46 \\
& \seqsplit{POST-запросов} (строка 26) & 47 \\
& \seqsplit{POST-запросов.} (строка 29) & 47 \\
& \seqsplit{(Postgres)Количество} (строка 5) & 49 \\
& \seqsplit{мсPostgres} (строка 28) & 49 \\
& \seqsplit{c-средних.} (строка 32) & 50 \\
& \seqsplit{c-средних} (строка 19) & 54 \\
& \seqsplit{("BBНеактивен} (строка 8) & 61 \\
& \seqsplit{("BBактивен} (строка 38) & 63 \\
& \seqsplit{("BBнеактивен} (строка 9) & 64 \\ \hline

% _09.pdf
\_09.pdf & \seqsplit{ми(xtarget,ytarget,zde} (строка 28) & 4 \\
& \seqsplit{деσ} (строка 8) & 5 \\
& \seqsplit{деα} (строка 12) & 5 \\
& \seqsplit{деP=} (строка 48) & 5 \\
& \seqsplit{ER-диаграмма} (строка 28) & 6 \\
& \seqsplit{ER-диаграмма} (строка 1) & 7 \\
& \seqsplit{CID-транзакций} (строка 26) & 18 \\
& \seqsplit{SQL-инъекций;} (строка 7) & 20 \\ \hline

% ВКР Селез.pdf
ВКР Селез.pdf & \seqsplit{XIV-ХХI} (строка 40) & 6 \\
& \seqsplit{Oб} (строка 27) & 7 \\
& \seqsplit{Вести.ru:} (строка 10) & 11 \\
& \seqsplit{Oб} (строка 24) & 30 \\
& \seqsplit{XIV-ХХI} (строка 29) & 31 \\
& \seqsplit{Вести.ru:} (строка 31) & 33 \\
& \seqsplit{Oб} (строка 13) & 68 \\
& \seqsplit{XIV-ХХI} (строка 29) & 70 \\
& \seqsplit{Вести.ru:} (строка 25) & 78 \\ \hline

% _03.pdf
\_03.pdf & \seqsplit{AиBразмерностей} (строка 5) & 3 \\
& \seqsplit{m×nиn×qсоответственно.} (строка 5) & 3 \\
& \seqsplit{Cразмерности} (строка 6) & 3 \\
& \seqsplit{Cкак} (строка 11) & 3 \\
& \seqsplit{i-й} (строка 11) & 3 \\
& \seqsplit{Aиj-го} (строка 11) & 3 \\
& \seqsplit{m·qэлементов} (строка 14) & 3 \\
& \seqsplit{nумножений} (строка 14) & 3 \\
& \seqsplit{(n−1)сложений.} (строка 14) & 3 \\
& \seqsplit{nскалярное} (строка 19) & 3 \\
& \seqsplit{n)иv=} (строка 20) & 3 \\
& \seqsplit{nвводят} (строка 2) & 4 \\
& \seqsplit{t=⌊n/2⌋и} (строка 2) & 4 \\
& \seqsplit{cij(введём} (строка 7) & 4 \\
& \seqsplit{еслиnнечётно} (строка 13) & 4 \\
& \seqsplit{еслиnчётно} (строка 14) & 4 \\
& \seqsplit{fif=fусловия} (строка 2) & 8 \\
& \seqsplit{fA,если} (строка 4) & 8 \\
& \seqsplit{fB,иначе.(2.3)} (строка 5) & 8 \\
& \seqsplit{ffor=fинициализации} (строка 7) & 8 \\
& \seqsplit{+fсравнения} (строка 7) & 8 \\
& \seqsplit{+N(fтела+fинкремента} (строка 7) & 8 \\
& \seqsplit{cреда} (строка 4) & 10 \\
& \seqsplit{иrandom} (строка 19) & 17 \\ \hline

% _02.pdf
\_02.pdf & \seqsplit{гдеX-} (строка 10) & 4 \\
& \seqsplit{i-е} (строка 11) & 4 \\
& \seqsplit{XиY:} (строка 18) & 4 \\
& \seqsplit{гдеµXиµY-} (строка 22) & 4 \\
& \seqsplit{XиYсоответственно.} (строка 22) & 4 \\
& \seqsplit{гдеσXиσY-} (строка 30) & 4 \\
& \seqsplit{XиY.} (строка 30) & 4 \\
& \seqsplit{гдеdi-} (строка 4) & 5 \\
& \seqsplit{XиY,n-} (строка 4) & 5 \\
& \seqsplit{топ-N} (строка 22) & 6 \\
& \seqsplit{meanмежду} (строка 30) & 23 \\ \hline

% _00.pdf
\_00.pdf & \seqsplit{input("Введите} (строка 35) & 7 \\ \hline

% _65-1.pdf
\_65-1.pdf & \seqsplit{иB} (строка 10) & 4 \\
& \seqsplit{иB} (строка 5) & 5 \\
& \seqsplit{деani,} (строка 27) & 5 \\
& \seqsplit{дoёмких} (строка 30) & 5 \\
& \seqsplit{деM,N,K} (строка 7) & 7 \\
& \seqsplit{дoёмк} (строка 8) & 11 \\
& \seqsplit{дoёмк} (строка 10) & 11 \\
& \seqsplit{дoёмк} (строка 16) & 11 \\
& \seqsplit{сравнения+N(f} (строка 17) & 11 \\
& \seqsplit{тела+f} (строка 17) & 11 \\
& \seqsplit{дoёмк} (строка 18) & 11 \\
& \seqsplit{дoёмк} (строка 11) & 12 \\
& \seqsplit{дoёмк} (строка 12) & 12 \\
& \seqsplit{дoёмк} (строка 13) & 12 \\
& \seqsplit{иCOLS} (строка 19) & 12 \\
& \seqsplit{дoёмк} (строка 19) & 12 \\
& \seqsplit{дoёмк} (строка 1) & 13 \\
& \seqsplit{дoёмк} (строка 4) & 13 \\
& \seqsplit{дoёмк} (строка 7) & 13 \\
& \seqsplit{дoёмк} (строка 11) & 13 \\
& \seqsplit{дoёмк} (строка 16) & 13 \\
& \seqsplit{дoёмк} (строка 1) & 14 \\
& \seqsplit{дoёмк} (строка 3) & 14 \\
& \seqsplit{дoёмк} (строка 6) & 14 \\
& \seqsplit{дoёмк} (строка 9) & 14 \\
& \seqsplit{деflast} (строка 12) & 14 \\
& \seqsplit{дoёмк} (строка 12) & 14 \\
& \seqsplit{дoёмк} (строка 17) & 14 \\
& \seqsplit{(N)Время} (строка 2) & 27 \\
& \seqsplit{дoёмких} (строка 5) & 29 \\ \hline

% _04.pdf
\_04.pdf & \seqsplit{Cценарий} (строка 22) & 3 \\
& \seqsplit{работеThe} (строка 2) & 7 \\
& \seqsplit{humor):Использование} (строка 6) & 7 \\
& \seqsplit{1)amusement(развлечение)} (строка 27) & 7 \\
& \seqsplit{2)enjoyment(наслаждение,} (строка 28) & 7 \\
& \seqsplit{3)fun(веселье,} (строка 29) & 7 \\
& \seqsplit{4)pleasure(удовольствие)} (строка 30) & 7 \\
& \seqsplit{5)absurdity(нелепость,} (строка 31) & 7 \\
& \seqsplit{6)irony(ирония)} (строка 32) & 7 \\
& \seqsplit{7)laughableness(забавность)} (строка 33) & 7 \\
& \seqsplit{8)ludicrousness(нелепость)} (строка 34) & 7 \\
& \seqsplit{9)ridiculousness(нелепость)} (строка 35) & 7 \\
& \seqsplit{10)whimsicality(прихотливость,} (строка 36) & 7 \\
& \seqsplit{11)wittiness(остроумие)} (строка 1) & 8 \\
& \seqsplit{12)wryness(искаженность,} (строка 2) & 8 \\
& \seqsplit{13)burlesque(бурлеск)} (строка 3) & 8 \\
& \seqsplit{14)caricature(карикатура)} (строка 4) & 8 \\
& \seqsplit{15)farce(фарс)} (строка 5) & 8 \\
& \seqsplit{16)jest(шутка,} (строка 6) & 8 \\
& \seqsplit{17)lampoon(памфлет,} (строка 7) & 8 \\
& \seqsplit{18)parody(пародия)} (строка 8) & 8 \\
& \seqsplit{19)satire(сатира)} (строка 9) & 8 \\
& \seqsplit{обработки:cистема} (строка 27) & 8 \\
& \seqsplit{ищображения:cистема} (строка 29) & 8 \\
& \seqsplit{AI-модель} (строка 21) & 10 \\
& \seqsplit{AI-алгоритмы} (строка 5) & 11 \\
& \seqsplit{AI-системы} (строка 12) & 12 \\
& \seqsplit{AI-генерируемый} (строка 20) & 12 \\
& \seqsplit{Python-скрипт} (строка 4) & 13 \\
& \seqsplit{Cкрипт} (строка 4) & 13 \\
& \seqsplit{GIFДа} (строка 15) & 16 \\
& \seqsplit{generatorГенерация} (строка 27) & 16 \\
& \seqsplit{(aigo.tools)Генерация} (строка 36) & 16 \\
& \seqsplit{MachineГенерация} (строка 9) & 17 \\
& \seqsplit{CaptionsГенерация} (строка 15) & 17 \\
& \seqsplit{CaptioningГенерация} (строка 25) & 17 \\
& \seqsplit{HumorГенерация} (строка 40) & 17 \\
& \seqsplit{AdaptabilityГенерация} (строка 49) & 17 \\
& \seqsplit{GeneratorГенерация} (строка 54) & 17 \\
& \seqsplit{cuptionsГенерация} (строка 59) & 17 \\
& \seqsplit{GeneratorГенерация} (строка 8) & 18 \\
& \seqsplit{IDEF0-диаграмма} (строка 5) & 22 \\
& \seqsplit{IDEF0-диаграмма} (строка 7) & 22 \\
& \seqsplit{xиG} (строка 25) & 24 \\
& \seqsplit{направлениеΘзатем} (строка 36) & 24 \\
& \seqsplit{aрктангенс,} (строка 43) & 24 \\
& \seqsplit{пикселяg} (строка 2) & 25 \\
& \seqsplit{cс} (строка 2) & 25 \\
& \seqsplit{соседейg} (строка 2) & 25 \\
& \seqsplit{pна} (строка 2) & 25 \\
& \seqsplit{гдеs(x)} (строка 6) & 25 \\
& \seqsplit{1еслиx≥0} (строка 6) & 25 \\
& \seqsplit{сивностиh(k)для} (строка 15) & 25 \\
& \seqsplit{гдеn} (строка 17) & 25 \\
& \seqsplit{ЕслиF} (строка 24) & 25 \\
& \seqsplit{c-й} (строка 24) & 25 \\
& \seqsplit{позиции(i,} (строка 24) & 25 \\
& \seqsplit{картыP} (строка 25) & 25 \\
& \seqsplit{cравен:} (строка 25) & 25 \\
& \seqsplit{гдеHиW—} (строка 30) & 25 \\
& \seqsplit{частиx} (строка 34) & 25 \\
& \seqsplit{iвычисляется} (строка 34) & 25 \\
& \seqsplit{вниманияα} (строка 34) & 25 \\
& \seqsplit{контекстC} (строка 34) & 25 \\
& \seqsplit{Весаα} (строка 1) & 26 \\
& \seqsplit{iвычисляются} (строка 1) & 26 \\
& \seqsplit{частьюx} (строка 2) & 26 \\
& \seqsplit{признакиx.} (строка 16) & 26 \\
& \seqsplit{классуC} (строка 17) & 26 \\
& \seqsplit{kпринимается} (строка 17) & 26 \\
& \seqsplit{вероятностиP(C} (строка 18) & 26 \\
& \seqsplit{zkдля} (строка 19) & 26 \\
& \seqsplit{классаkпреобразуются} (строка 19) & 26 \\
& \seqsplit{вx,} (строка 24) & 26 \\
& \seqsplit{метокs=} (строка 27) & 26 \\
& \seqsplit{N)для} (строка 27) & 26 \\
& \seqsplit{вероятностьP(s|x)или,} (строка 28) & 26 \\
& \seqsplit{энергиюE(s,x):} (строка 29) & 26 \\
& \seqsplit{i-го} (строка 10) & 27 \\
& \seqsplit{последовательностиSна} (строка 9) & 28 \\
& \seqsplit{кенамиt} (строка 10) & 28 \\
& \seqsplit{гдеc} (строка 3) & 29 \\
& \seqsplit{j-й} (строка 3) & 29 \\
& \seqsplit{токенt} (строка 6) & 29 \\
& \seqsplit{iявляется} (строка 6) & 29 \\
& \seqsplit{словаряV,V=\{v} (строка 6) & 29 \\
& \seqsplit{аN–} (строка 6) & 29 \\
& \seqsplit{(обычноN≤L).} (строка 7) & 29 \\
& \seqsplit{токенуt} (строка 8) & 29 \\
& \seqsplit{iиз} (строка 8) & 29 \\
& \seqsplit{словаряVсопоставляется} (строка 8) & 29 \\
& \seqsplit{ПустьT} (строка 15) & 29 \\
& \seqsplit{токенуv} (строка 17) & 29 \\
& \seqsplit{jиз} (строка 17) & 29 \\
& \seqsplit{словаряV(размером|V|)} (строка 17) & 29 \\
& \seqsplit{гдеd–} (строка 19) & 29 \\
& \seqsplit{строкаW} (строка 24) & 29 \\
& \seqsplit{e,jэтой} (строка 24) & 29 \\
& \seqsplit{токеновT} (строка 1) & 30 \\
& \seqsplit{IDпреобразуется} (строка 1) & 30 \\
& \seqsplit{гдеe} (строка 4) & 30 \\
& \seqsplit{эмбеддинговX=} (строка 5) & 30 \\
& \seqsplit{черезLидентичных} (строка 6) & 30 \\
& \seqsplit{токенx} (строка 7) & 30 \\
& \seqsplit{iформирует} (строка 7) & 30 \\
& \seqsplit{Запросq} (строка 7) & 30 \\
& \seqsplit{Ключk} (строка 7) & 30 \\
& \seqsplit{Значениеv} (строка 8) & 30 \\
& \seqsplit{Выходz} (строка 8) & 30 \\
& \seqsplit{iдля} (строка 8) & 30 \\
& \seqsplit{токенаi(одна} (строка 8) & 30 \\
& \seqsplit{гдеQ,K,V–} (строка 12) & 30 \\
& \seqsplit{Attentionhтаких} (строка 12) & 30 \\
& \seqsplit{проекциямиWQ} (строка 13) & 30 \\
& \seqsplit{FFN:Применяется} (строка 19) & 30 \\
& \seqsplit{ходуz} (строка 20) & 30 \\
& \seqsplit{iслоя} (строка 20) & 30 \\
& \seqsplit{соединения(x+Sublayer(x))} (строка 22) & 30 \\
& \seqsplit{слоям(Layer} (строка 23) & 30 \\
& \seqsplit{последнегоL-го} (строка 23) & 30 \\
& \seqsplit{векторовH=} (строка 24) & 30 \\
& \seqsplit{последовательностиY=} (строка 27) & 30 \\
& \seqsplit{M)токен} (строка 27) & 30 \\
& \seqsplit{токенаy} (строка 28) & 30 \\
& \seqsplit{tна} (строка 28) & 30 \\
& \seqsplit{шагеt.} (строка 28) & 30 \\
& \seqsplit{промтаHи} (строка 30) & 30 \\
& \seqsplit{логитовl} (строка 31) & 30 \\
& \seqsplit{t∈R|V|для} (строка 31) & 30 \\
& \seqsplit{вероятностейP(y} (строка 32) & 30 \\
& \seqsplit{t|y<t,H)с} (строка 32) & 30 \\
& \seqsplit{гдеl} (строка 37) & 30 \\
& \seqsplit{дляj-го} (строка 37) & 30 \\
& \seqsplit{шагеt.} (строка 37) & 30 \\
& \seqsplit{токенаy} (строка 38) & 30 \\
& \seqsplit{t:Из} (строка 38) & 30 \\
& \seqsplit{распределенияP(y} (строка 38) & 30 \\
& \seqsplit{t|y<t,H)выбира-} (строка 38) & 30 \\
& \seqsplit{токенy} (строка 2) & 31 \\
& \seqsplit{tдобавляется} (строка 2) & 31 \\
& \seqsplit{генерацииy} (строка 3) & 31 \\
& \seqsplit{Cценарий} (строка 6) & 31 \\
& \seqsplit{IDEF0-диаграммы} (строка 14) & 31 \\
& \seqsplit{функцииget\_image\_caption.} (строка 26) & 35 \\
& \seqsplit{функцииgenerate\_funny\_caption\_llm,} (строка 24) & 37 \\
& \seqsplit{1)mistral-7b...} (строка 20) & 40 \\
& \seqsplit{2)meta-llama...} (строка 22) & 40 \\
& \seqsplit{женияMistral} (строка 4) & 41 \\
& \seqsplit{Время(Cек)} (строка 4) & 41 \\
& \seqsplit{VRAM(Гб)} (строка 4) & 41 \\
& \seqsplit{мя(Сек)Llama} (строка 6) & 41 \\
& \seqsplit{Clause)Ограничено} (строка 24) & 50 \\
& \seqsplit{"Python"и} (строка 10) & 51 \\
& \seqsplit{Моcква,} (строка 8) & 62 \\ \hline

% _10.pdf
\_10.pdf & \seqsplit{помощьюDefinite} (строка 9) & 3 \\
& \seqsplit{nи} (строка 20) & 3 \\
& \seqsplit{всеB} (строка 21) & 3 \\
& \seqsplit{приn=} (строка 21) & 3 \\
& \seqsplit{фактыresult\_type/4.} (строка 26) & 3 \\
& \seqsplit{B—примертерма} (строка 2) & 4 \\
& \seqsplit{B=Aθдля} (строка 2) & 4 \\
& \seqsplit{Number,Gender,Personконкретизируются} (строка 4) & 4 \\
& \seqsplit{пример,noun/4иverb/4унифицируются} (строка 9) & 4 \\
& \seqsplit{DCG-правила} (строка 13) & 4 \\
& \seqsplit{intдаютfloat,} (строка 18) & 4 \\
& \seqsplit{какint.} (строка 19) & 4 \\
& \seqsplit{1object(автобус,} (строка 21) & 4 \\
& \seqsplit{2object(брат,} (строка 22) & 4 \\
& \seqsplit{3object(вода,} (строка 23) & 4 \\
& \seqsplit{4object(вопрос,} (строка 24) & 4 \\
& \seqsplit{DCG-правило} (строка 26) & 4 \\
& \seqsplit{7noun(брат,} (строка 27) & 4 \\
& \seqsplit{object(брат,} (строка 29) & 4 \\
& \seqsplit{noun/4,verb/4иobject/5} (строка 6) & 5 \\
& \seqsplit{(DCG)).Грамматика} (строка 8) & 5 \\
& \seqsplit{Personпри} (строка 18) & 5 \\
& \seqsplit{переменныеNumber,} (строка 27) & 5 \\
& \seqsplit{Personпрохо-} (строка 27) & 5 \\
& \seqsplit{типизации:рекурсивнаяDCG,} (строка 3) & 6 \\
& \seqsplit{термовnoun/4иverb/4по} (строка 7) & 7 \\
& \seqsplit{floatи} (строка 13) & 7 \\
& \seqsplit{обобщающийnum.} (строка 13) & 7 \\
& \seqsplit{фактамиresult\_type/4.} (строка 21) & 7 \\
& \seqsplit{предикатаresult\_type/4} (строка 3) & 8 \\
& \seqsplit{подставляетint,} (строка 6) & 8 \\
& \seqsplit{1noun(брат,} (строка 7) & 9 \\
& \seqsplit{2noun(вопросы,} (строка 8) & 9 \\
& \seqsplit{4verb(работал,} (строка 10) & 9 \\
& \seqsplit{5verb(решается,} (строка 11) & 9 \\
& \seqsplit{DCG-правило} (строка 21) & 9 \\
& \seqsplit{DCG–правила} (строка 24) & 9 \\
& \seqsplit{какfloat.} (строка 26) & 9 \\
& \seqsplit{типintпо} (строка 27) & 9 \\
& \seqsplit{noun(брат,} (строка 4) & 11 \\
& \seqsplit{иverb(работал,} (строка 5) & 11 \\
& \seqsplit{→true:noun(братья,} (строка 7) & 11 \\
& \seqsplit{иverb(работают,} (строка 7) & 11 \\
& \seqsplit{third)унифицируются.} (строка 8) & 11 \\
& \seqsplit{:noun(вопросы,} (строка 9) & 11 \\
& \seqsplit{иverb(решается,} (строка 9) & 11 \\
& \seqsplit{present)конфликтуют} (строка 10) & 11 \\
& \seqsplit{noun(брат,} (строка 14) & 11 \\
& \seqsplit{братverb(Lemma,} (строка 15) & 11 \\
& \seqsplit{verb(работал,} (строка 16) & 11 \\
& \seqsplit{5→int(еслиxнеизвестен,intпо} (строка 1) & 12 \\
& \seqsplit{+\_→int(эвристика} (строка 5) & 12 \\
& \seqsplit{pluralиsingular} (строка 17) & 12 \\
& \seqsplit{системаможетбыть...} (строка 15) & 13 \\
& \seqsplit{Особоевнимание...} (строка 7) & 15 \\
& \seqsplit{DCG–синтаксис} (строка 3) & 16 \\ \hline

% _06.pdf
\_06.pdf & \seqsplit{IDEFО} (строка 10) & 1 \\
& \seqsplit{чaс} (строка 8) & 2 \\
& \seqsplit{IDEFО} (строка 25) & 2 \\
& \seqsplit{IDEFО} (строка 7) & 10 \\ \hline

% _07.pdf
\_07.pdf & \seqsplit{программированияPython[1]} (строка 3) & 7 \\
& \seqsplit{классPoliticalRiskAnalyzer,} (строка 20) & 7 \\ \hline

\end{longtable}
\endgroup
